% Transcription — fonds Alexandre Grothendieck, shelfmark 67, batch 1 (pages 1-20).
% Established from the facsimile archives/batches/67/batch-01.pdf (a mirror of
% https://grothendieck.umontpellier.fr/67.pdf), per the skill
% transcribe-grothendieck. The batch file carries no cover sheet: PDF page N is
% archivists' page N; the pencilled numbers bottom-left ("2", "3", "4", "6",
% "10", "12", "15" ... "20") read at the PDF page of the same number.
%
% Pass: Opus 5.5 (claude-opus-5-5), 2026-09-22 — first pass, unchecked against the pages by a human.
%
% What the batch is.
%   1      Cover in his hand: « Chirurgie des surfaces conformes ... (1978 ou
%          77 ?) ». No \page; its words open the section.
%   2-4, 6 A numbered programme, 1) to 7), for surgery on conformal
%          surfaces: découpage along a rectifiable 1-complex, recollage along
%          an admissible equivalence relation on the boundary, rétrécissage
%          and élargissage (removing / adding a collar), a mixed operation,
%          then 7) moduli of S ⊂ K ⊂ X rigidified against a standard
%          situation. Page 6 continues page 4 (4 ends on « Il », 6 opens on
%          « vaut mieux ») and counts real moduli, N - (3g-3) as a sum over
%          the faces f of X - K. Fast hand; the formulas carry, much of the
%          connecting prose does not.
%   5      Sketch sheet, scanned upside down: lim over K - L of K_x, the
%          covering K~ -> K, Σ(X_i, ∂X_i) = |K~_i|, drawings of pants and an
%          orange-sliced sphere.
%   7      Unrelated typed verso (a list of plots of land). Absent.
%   8-13   Scratch on envelopes: three successive drafts (8, 10, 12) of one
%          list of six equivalent data — elliptic curve with oriented closed
%          geodesic; real torus Γ_0; primitive class in π_1; forms of G_m and
%          discrete Z ⊂ C*; projective line with a non-parabolic u; forms of
%          GL(1) with a discrete non-unipotent Z — and the chain of
%          bijections {Γ} ≃ ... linking them. 9: a sketch of discs around
%          points. 11: envelope front, no mathematics, absent. 13: the front
%          of 12's envelope, postmarked Nagoya 20.12.77, with computations on
%          λ + λ^{-1}.
%   14     Blank. Absent.
%   15-20  Typed letter (carbon or copy, unsigned on the page), « Les
%          Aumettes 15.4.1984 », « Dear Lipman Bers »: Teichmüller space
%          M~_{g,ν} against Thurston's MB~_{g,ν} of surfaces with boundary,
%          generalized boundary, admissible systems of discs, radii, the model
%          of three orange slices on the Riemann sphere for pants, and the
%          conjecture MB~ ≃ M~ × [0,1[^I. In English; his handwritten
%          corrections are kept, the Greek letters (ν, λ, ρ, l) are inked by
%          hand in the typescript.
\documentclass[11pt,a4paper]{article}
% Le préambule commun à toutes les transcriptions du fonds Grothendieck.
%
% Il tient en une idée : **l'incertitude se compose, elle ne se résout pas.**
% Une transcription qui lisse un mot douteux a détruit la seule chose pour
% laquelle on la faisait — un lecteur doit pouvoir savoir, à chaque endroit,
% ce qui était sur le feuillet et ce qui a été ajouté. Les six macros qui
% suivent sont donc l'appareil critique, et non de la décoration.
%
% Les mêmes commandes sont reconnues par `scripts/render.mjs`, qui produit la
% vue de lecture du site. Une macro ajoutée ici doit l'être aussi là-bas,
% faute de quoi le rendu échouera bruyamment — ce qui est voulu : un rendu
% silencieusement faux est pire qu'un rendu absent.

\ProvidesPackage{grothendieck}[2026/08/08 Transcriptions du fonds Grothendieck]

\usepackage{amsmath,amssymb,amsthm}
\usepackage{mathrsfs}
\usepackage{stmaryrd}
% Les diagrammes commutatifs sont partout dans ce fonds : c'est la langue
% même de la géométrie algébrique de Grothendieck, et une transcription qui
% les rendrait en prose ne transcrirait rien.
\usepackage{tikz-cd}
% Français par défaut — c'est la langue des feuillets — mais l'anglais est
% chargé aussi : la traduction et le résumé sont des documents anglais, et la
% typographie française (espaces avant les deux-points) y serait une faute.
% Ces documents basculent par \selectlanguage{english} en tête de corps.
\usepackage[english,french]{babel}
\usepackage{fontspec}
\usepackage{geometry}
\geometry{a4paper,margin=25mm,marginparwidth=28mm}
\usepackage{marginnote}
\usepackage{xcolor}
% ulem, not soul. `soul`'s \st and \ul are built on a character-by-character
% scan that XeLaTeX's font machinery breaks: the document fails with an
% undefined control sequence rather than degrading. ulem does the same two jobs
% and is Unicode-engine safe. `normalem` stops it hijacking \emph, which the
% transcriptions use for Grothendieck's own emphasis.
\usepackage[normalem]{ulem}
\usepackage{hyperref}
\hypersetup{colorlinks=true,linkcolor=black,urlcolor=black,pdfborder={0 0 0}}

% --- Métadonnées du lot -----------------------------------------------------
% Reprises telles quelles par le script de rendu, qui les affiche en tête de
% la vue de lecture. Elles fixent ce que le fichier prétend couvrir : sans
% elles, une transcription flotte sans point d'attache dans un fonds de seize
% mille feuillets.

\newcommand{\folder}[1]{\gdef\grothFolder{#1}}
\newcommand{\batch}[1]{\gdef\grothBatch{#1}}
\newcommand{\leaves}[2]{\gdef\grothFirst{#1}\gdef\grothLast{#2}}
\newcommand{\foldertitle}[1]{\gdef\grothFoldertitle{#1}}
\gdef\grothFolder{?}\gdef\grothBatch{?}\gdef\grothFirst{?}\gdef\grothLast{?}\gdef\grothFoldertitle{}

% --- Le résumé --------------------------------------------------------------
%
% Il ouvre la lecture modernisée, sous le titre « Résumé », et il est composé
% en petit. Ce n'est pas un ornement : c'est le seul passage du document qui
% ne soit pas des mathématiques, et un lecteur doit pouvoir le reconnaître
% avant de l'avoir lu — puis le sauter s'il connaît déjà le sujet, ou s'y
% arrêter s'il ne le connaît pas. Le filet qui le suit marque la reprise du
% corps.

\newenvironment{resume}
  {\small\setlength{\parskip}{0.45em}}
  {\par\medskip\hrule\medskip}

% --- Mots-clés ---------------------------------------------------------------
%
% En anglais, en clôture du résumé de la lecture modernisée : le vocabulaire
% moderne sous lequel chercher ce que les feuillets construisent. Le manifeste
% du site (`npm run manifest`) les extrait pour étiqueter la cote entière —
% cette ligne est la source unique des tags, il n'y a pas de fichier à côté.
\newcommand{\keywords}[1]{\par\smallskip\noindent
  {\footnotesize\textsc{Keywords} --- \textit{#1}}\par}

% --- L'appareil critique ----------------------------------------------------

% Le numéro de feuillet, en marge. C'est l'ancre qui permet de régler un
% désaccord sur une formule en regardant *un* feuillet plutôt qu'une chemise
% de six cents. Le site s'en sert aussi pour tourner les pages du fac-similé
% quand on fait défiler la transcription.
\newcommand{\leaf}[1]{%
  \marginnote{\footnotesize\color{gray!70}\textsf{#1}}%
  \label{leaf:#1}\ignorespaces}

% Illisible : on ne devine pas. Un mot inventé qui se lit comme les autres est
% le pire résultat possible.
\newcommand{\ill}{\textcolor{red!60!black}{[\,\ldots\,]}}

% Lecture douteuse : le mot est donné, et signalé comme incertain.
\newcommand{\uncertain}[1]{\uline{#1}}

% Ajout de l'éditeur — une lettre manquante, un mot sous-entendu.
\newcommand{\add}[1]{\textcolor{blue!50!black}{[#1]}}

% Biffé par Grothendieck lui-même. Ce qu'il a rayé fait partie du document :
% une rature dit souvent où la pensée a changé de direction.
\newcommand{\struck}[1]{\sout{#1}}

% Note du transcripteur, sur l'état matériel du feuillet.
\newcommand{\note}[1]{\par\smallskip{\small\color{gray!60!black}\textit{[#1]}}\par\smallskip}

% Note marginale de Grothendieck, distincte du corps du texte.
\newcommand{\marginal}[1]{\par\smallskip\noindent\hspace{1em}%
  {\small\color{gray!60!black}#1}\par\smallskip}

% --- Titre ------------------------------------------------------------------

\AtBeginDocument{%
  \begin{center}
    {\large\bfseries Fonds Alexandre Grothendieck --- cote n\textsuperscript{o}~\grothFolder}\\[2pt]
    {\itshape\grothFoldertitle}\\[4pt]
    {\small feuillets \grothFirst--\grothLast\ (lot \grothBatch)}\\[2pt]
    {\footnotesize\color{gray!60!black}Transcription établie d'après le fac-similé numérisé
     par l'Université de Montpellier.\\
     Les lectures incertaines sont soulignées, les illisibles notées [\,\ldots\,],
     les ajouts entre crochets.}
  \end{center}
  \vspace{1em}}

\endinput


\folder{67}
\batch{1}
\pages{1}{20}
\dating{[à partir de 1977]-1984}
\watermark{Édition de démonstration}
\foldertitle{Chirurgie des surfaces conformes : notes manuscrites (s.d., 1983-1984), lettres (1984)}

\begin{document}

\section{Chirurgie des surfaces conformes}
\note{les mots du titre sont de sa main, en haut à droite du feuillet 1, une
chemise par ailleurs vierge qui ne reçoit pas de numéro de page ; le titre y
est suivi de « \ill{} (1978 ou 77 ?) », de sa main aussi}

\page{2}
Chirurgie des surfaces conformes.

1) \emph{Découpage} : Découpage d'une surface conforme \struck{\ill{}}
(\uncertain{évent.} avec bord) suivant un sous-1-complexe rectifiable
compact : On trouve une surface \struck{conforme} \uncertain{à} bord, avec un
\uncertain{morceau} \uncertain{compact} de bord compact marqué.

2) \emph{Recollage} : À partir d'une surface conforme (avec bord
\struck{\ill{}}) \uncertain{et} d'un morceau compact du bord, et une relation
d'équivalence “admissible” \uncertain{(au sens rectifiable)} sur le bord,
\struck{\ill{}} on trouve une surface conforme avec bord…

3) \emph{Rétrécissage}. $K \subset X$ (1-compl.\ cpct, s.\ conforme avec
bord)

$\widetilde{X}$ découpé de $X$ suivant $K$, $\partial'\widetilde{X}$ partie
distinguée du bord, on choisit un \uncertain{multi-}anneau autour de
$\partial'\widetilde{X}$, \struck{et alors} (un germe de tel anneau autour
d'une composante $(\partial'_{i}\widetilde{X})_{i}$ revient : la donnée d'une
structure \struck{riemannienne} \struck{\uncertain{régulière}} rectifiable de
long.\ totale 1 sur celle-ci….) $\widetilde{X}^{*}$ la surface
\uncertain{rétrécie} \struck{\ill{}} \uncertain{à bord} \ill{} déduite
\ill{} la couronne.

Choisissons un isomorphisme rectifiable
$\partial'\widetilde{X}^{*} \xrightarrow{\ \sim\ } \partial'\widetilde{X}$
compatible avec les composantes, et utilisons la relation d'équivalence
\uncertain{de recollement} sur $\partial'\widetilde{X}$ pour en déduire une
sur $\widetilde{X}^{*}$. On en conclut \struck{\ill{}} $X^{*}$ et
$K \subset X^{*}$.

\note{à gauche, un croquis de $\widetilde{X}$ : deux anneaux hachurés le long
de deux composantes du bord, l'épaisseur de chacun marquée $\varepsilon_{i}$,
et sous le dessin « $0 \leqslant \varepsilon_{i} < +\infty$ »}

\page{3}
4) \emph{Élargissage}. Opération inverse — au lieu de retrancher une
couronne, on en rajoute une…

5) Opération mixte : \emph{\ill{}}
\note{le nom de l'opération, souligné, commence par un p et finit en -age ; il
revient au 4') et au 6) ci-dessous et n'a pu être lu}

On considère un voisinage tubulaire \uncertain{standard} \ill{} de
$\partial'\widetilde{X}$, et dans celui-ci, un cercle $\Gamma_{i}$, et un
isom.\ rect.\ $\Gamma_{i} \simeq (\partial'\widetilde{X})_{i}$. On
\struck{\ill{}} \uncertain{définit} \struck{\ill{}} \uncertain{considère}
$\widetilde{X}^{*}$ déduit de $\widetilde{X}$ en enlevant \uncertain{$K_{i}$}
“la \struck{\ill{}} partie de \uncertain{$T_{i}$} entourée par $\Gamma_{i}$”.
On termine comme dans 3°), 4°) pour former $\widetilde{X}^{*}$ et
$K \subset X^{*}$.

\marginal{C'est une composition rétrécissage $\circ$ élargissage
\uncertain{ou} élargissage $\circ$ rétrécissage, \ill{} au choix}

\note{à gauche, un croquis : une composante $(\partial'\widetilde{X})_{i}$
entourée d'une couronne hachurée, et dans celle-ci le cercle $\Gamma_{i}$}

4') \emph{Commutativité} de l'opération \ill{} relativement à diverses
composantes connexes de $K$.

6) Cas d'une composante \uncertain{de $K$ bordée} \uncertain{plongée
$\mathbb{R}$-analytique} isolée. L'opération de \ill{} \uncertain{revient}
donc à une composition de 1°) \ill{} \uncertain{une couronne intérieure}
\uncertain{autour de} $C$, 2°) \uncertain{$\mathbb{R}$-analytiques}
\uncertain{Recoller sur} un cylindre les deux \uncertain{bords}
\uncertain{joints} de la couronne \uncertain{enlevée}…

\note{à gauche, trois croquis de cylindres traversés de cercles et de
courbes, le troisième biffé}

\uncertain{Relations} avec les \struck{\ill{}} \ill{} plongés \uncertain{en
forme} d'une petite couronne…

\page{4}
7) Modules pour $S \subset K \subset X$

On suppose \uncertain{compact, et} pour simplifier $X$ \struck{compact}
\uncertain{sans bord} (orientée).

Ceci dit, on fixe $S_{0} \subset K_{0} \subset X_{0}$ standard
\uncertain{top.}, et on \emph{rigidifie} $(S \subset K \subset X)$ en se
donnant une classe d'isotopie d'homéo.\ top.\ avec la situation standard. On
élimine ainsi les \uncertain{automorphismes}, sauf dans le cas
\struck{\ill{}} où $X_{0}$ a des composantes de \uncertain{genre} $0, 1$ ….

Pour avoir un “espace de modules”, on \uncertain{pourrait} exiger
\uncertain{v.ca.} que les arêtes de $K$ soient des arcs de cercle pour la
\uncertain{connexion} projective canonique. Mais les “cartes \ill{}”
\uncertain{ne} \uncertain{satisfont} pas à cette condition ! On pourrait
exiger plutôt que \ill{} \struck{\ill{}} les \uncertain{paramétrisations}
géodésiques sur le \uncertain{disque}, \ill{} disque avec l'\ill{}
\uncertain{hyperbolique} \ill{} de \uncertain{vecteurs} aux \ill{}.
[\uncertain{Éventuellement} on pourrait exiger l'équidistance \uncertain{des
sommets} \ill{} des bords de chaque composante \ill{} le disque…] Mais on a
des choix \uncertain{de disques} \ill{} par les morceaux \ill{} : pas de
\ill{} \uncertain{constant}. Il

\page{5}
\note{feuillet de croquis, formules éparses et dessins dans tous les sens ; on
n'en donne que les formules}

\struck{$\widetilde{L}_{0} = \pi^{-1}(L_{0}) \subset \widetilde{K}_{0}$}
\qquad \uncertain{$\widetilde{L} = T$}

\[
  \bigl|\widetilde{K}\bigr| - \bigl|\widetilde{L}\bigr| \xrightarrow{\ \sim\ } |K| - |L| ,
  \qquad
  (\widetilde{K}, \widetilde{L}) \rightsquigarrow
\]
“\ill{}”
\note{le mot entre guillemets suit la flèche $\rightsquigarrow$ sur la même
ligne}

\[
  \widetilde{K} \xrightarrow{\ \pi\ } K ,
  \qquad
  \widetilde{L} = \pi^{-1}(L) \longrightarrow L ,
  \qquad
  \widetilde{L} \subset \widetilde{K}
\]

\[
  \widetilde{K} = \varinjlim_{x \in K \setminus L} K_{x} ,
  \qquad
  K_{x} = \{\, y \in K \mid y \leqslant x \,\}
\]

\[
  \Sigma(X_{i}, \partial X_{i}) = \bigl|\widetilde{K}_{i}\bigr| ,
  \qquad \widetilde{K}_{i}
\]

\note{les dessins : une sphère découpée en fuseaux, des arcs emboîtés tracés
dans un disque, une rangée de disques accolés, un cercle muni d'une anse et un
embranchement en Y, c'est-à-dire des « pantalons » ; en marge, des mots
griffonnés illisibles et des coordonnées $(x', y')$}

\page{6}
vaut mieux dès lors exiger que ces \uncertain{arêtes} soient
\emph{\uncertain{horocycliques}} par morceaux….

Dans le premier cas (\struck{\ill{}} \uncertain{exigé} que les disques
\uncertain{ou} le disque \ill{} \ill{}…)

\ill{}, en \uncertain{supposant} $\pi_{0}(K_{0}) \to \pi_{0}(X_{0})$
surjectif, \struck{\ill{} \ill{} \ill{} des \ill{} \ill{}….} nb de modules
réels
\[
  N = \Bigl( \sum_{\substack{f \in \pi_{0}(X \setminus K) = F \\ f \ \text{pas un disque}}}
  \bigl( 6g(f) + 4\nu(f) - 6 \bigr) \Bigr)
\]
\note{un terme ajouté à droite de la parenthèse est biffé ; en dessous, une
ligne entière est biffée en zigzag, dont on lit seulement « \ill{} des $f$ qui
\ill{} des disques »}

\marginal{On \uncertain{suppose} $X_{0}$ \uncertain{connexe}, \ill{} $K$
\uncertain{connexe}, $\lambda$ \ill{} de $K$}

\uncertain{$\forall$} \ill{} $g \geqslant 2$ : $3g - 3$ où
\[
  2 - 2g = \chi = \sum_{f \in F} \bigl( 2 - 2g(f) - \nu(f) \bigr) + 1 - \lambda
\]
i.e.
\[
  2g - 2 = \sum_{f \in F} \bigl( 2g(f) - 2 + \nu(f) \bigr) + (\lambda - 1)
\]
\[
  g - 1 = \frac{1}{2} \sum_{f \in F} \bigl[ 2g(f) - 2 + \nu(f) \bigr] + \frac{3}{2}(\lambda - 1)
\]
\note{le $\tfrac{3}{2}$ de cette ligne paraît récrit sur un autre chiffre ;
la ligne précédente donnerait $\tfrac{1}{2}(\lambda-1)$}
\[
  3(g - 1) = \frac{3}{2} \sum_{f \in F} \bigl[ 2g(f) - 2 + \nu(f) \bigr] + \frac{3}{2}(\lambda - 1)
  = \sum_{f \in F} \Bigl[ 3\bigl(g(f) - 1\bigr) + \frac{3}{2}\nu(f) \Bigr] + \frac{3}{2}(\lambda - 1)
\]

$(g \geqslant 2)$
\[
  N - (3g - 3) = \sum_{\substack{f \in F \\ f \ \text{pas un disque}}}
  \Bigl( 3g(f) + \frac{5}{2}\nu(f) - 3 \Bigr) + \frac{3}{2}(1 - \lambda)
  \ \overset{?}{\geqslant}\ 0
\]
\note{une somme « $\sum_{f \in F,\, f \text{ disque}}$ » devant le dernier
terme est biffée ; dans $(1 - \lambda)$ le signe est récrit}

\ill{} \uncertain{des} $f$ qui \uncertain{sont} des disques

$(g = 1)$
\[
  N - 1 = \sum_{f \ \text{pas un disque}} \bigl( \underbrace{4\nu(f) - 6}_{\geqslant 2} \bigr) - 1
\]
[NB $g(f) = 0$ tjs]

\page{8}
\note{brouillon au dos d'une enveloppe ; première des trois versions de la
même liste, reprise aux pages 10 et 12}

(1) courbe ell.\ $X$ + géod.\ fermée \emph{orientée} $\Gamma$

(2) \struck{courbe} \uncertain{groupe} elliptique $X_{0}$ +
\struck{\uncertain{sous-tore réel}} $\mathbb{U}$ \ill{} + $\mathbb{U}$-torseur
$\Gamma$ \marginal{$\subset X_{0}(\mathbb{C})$}

(3) $X_{0}$ + $a \in \pi_{1}(X_{0})$ \uncertain{primitif} + $\Gamma$

(4)

Torseurs $\Gamma$ \uncertain{sous} $\mathbb{U}$, et $Z \subset
\mathbb{C}^{*}$ discret, $Z \simeq \mathbb{Z}$

(5)

(6)

\page{9}
\note{sur une enveloppe, un croquis sans légende : un grand disque
contenant un disque plus petit autour d'un point marqué, et deux petits
disques cerclés autour de deux autres points ; à côté, un second contour fermé
avec un point}

\page{10}
\note{deuxième version de la liste de la page 8}

(1) courbe ellipt.\ $X$ + géod.\ fermée orientée $\Gamma$
\uncertain{mod.\ tr.}

(2) —— $X$ + \struck{\ill{}} \ill{} $\mathbb{U}$ (tore compact standard)
$\subset X_{0}$

(3) —— $X$ + $a \in \pi_{1}(X)$ \uncertain{primitif}

(4) $Z \subset \mathbb{C}^{*}$ ($\Longleftrightarrow b \in \mathbb{C}^{*}$,
$0 < |b| < 1$) s.-gpe discret $(\simeq \mathbb{Z})$
\struck{$Z \subset \mathbb{C}^{*}$} et torseurs $\mathbb{T}$ sous
$\mathbb{C}^{*}$ \struck{de $Z \subset \mathbb{C}^{*}$ discret}

(5) Droite \uncertain{projective} \ill{} $E$, \struck{(NB
$\Delta = E^{*}$)}, et \struck{\ill{}} $b \in \mathbb{C}$, $0 < |b| < 1$

NB $\Delta = E^{*}$, $T = \mathbb{C}^{*}_{m}$, $Z$ = groupe engendré par $b$)

$P = \widehat{E}$, $\{u, u^{-1}\}$ = \struck{$b, b^{-1}$} $\{$homothéties
$b, b^{-1}\}$

(6) \struck{Formes $G$ de $\mathrm{Gl}(1)_{\mathbb{C}}$}, +
\struck{\ill{}} \struck{Borel $T \subset B \subset G$, \ill{}
$\in B$ \ill{}} \uncertain{(i.e.\ $G$ déployé)} \ill{} \uncertain{Couples de
Killing} \struck{le groupe \ill{}} de $\mathrm{Gl}(1)_{\mathbb{C}}$ ou
torseurs \uncertain{sous} $T \simeq \mathbb{C}^{*}$ \uncertain{(i.e.\ formes
de $\mathbb{C}^{*}$)} [\ill{} \ill{} \ill{} $E'$] \ill{}
$Z \subset T(\mathbb{C}) = \mathbb{C}^{*}$ \ill{} \ill{} $G \supset B \ni b$,
\ill{} \ill{} ½ 2-\ill{} (i.e.\ $Z$ \ill{} \ill{} $b$) et $|b|$
(défini \ill{} \ill{} : l'inverse) $\neq 1$
\note{le 6) est un palimpseste ; la version lisible est celle de la page 12}

\struck{\ill{}}

\[
  \text{donc } \{\Gamma\} \simeq X/\Gamma_{0} \simeq X/\mathbb{R}
  \simeq \Delta / T_{\mathbb{C}} \cdot Z
  \simeq \text{cercles de } P
\]
\note{sous les termes de la chaîne, les numéros cerclés (1), (2), (3'), (4),
(5) des données correspondantes ; un $X/\mathbb{Z}$ est corrigé en
$X/\mathbb{R}$, et après « cercles de $P$ » viennent, en petit :
« \struck{droites projectives $P$} \uncertain{admettant}
$\mathrm{Fix}(u, u^{-1})$ comme \ill{} bicentre »}

\page{12}
\note{troisième version de la liste, au dos d'une enveloppe (voir la page
13) ; les flèches doubles entre les numéros sont les siennes. En tête, une
ligne biffée : \struck{\ill{}}}

(1) courbe elliptique $X$ + géodésique fermée $\Gamma$ mod.\ translations

$\Updownarrow$

(2) courbe elliptique $X$ + \struck{géodésique fermée} \uncertain{s.-gpe
fermé} : i.e.\ tore réel $\Gamma_{0}$ de la courbe elliptique mod.\
translations

$\Updownarrow$

(3) courbe elliptique $X$ + \struck{\ill{}} \uncertain{sous-groupe}
\uncertain{facteur direct} \uncertain{libre de rg 1} de $\pi_{1}(X)$

$\Updownarrow$

(4) formes $T$ de $\mathbb{G}_{m}$, ($Z \subset T(\mathbb{C})$
\struck{\ill{}}, et torseurs sous $T$
\marginal{$b \in \mathbb{C}$, $0 < |b| < 1$}

(4') formes $T$ de $\mathbb{G}_{m}$, \struck{\ill{}
$\in \mathbb{C}$ \ill{}} et torseurs \uncertain{sous} $T$
\uncertain{complétés} $\{u, u^{-1}\}$ \ill{}

$\Updownarrow$

(5) droites projectives $P$, + \uncertain{automorph.} $u$ de $P$ non
paraboliques (i.e.\ ayant 2 pts fixes distincts) et ayant des val.\ propres
$\lambda, \lambda^{-1}$ telles que $|\lambda| \neq 1$ \uncertain{CQFS},
\ill{} : \struck{\ill{}} $\lambda$ \ill{} pour $|\lambda| > 1$)

$\Updownarrow$

(6) formes $G$ de $\mathrm{Gl}(1)_{\mathbb{C}}$, + s.-groupe discret
$(\simeq \mathbb{Z})$ $Z \subset G(\mathbb{C})$ \struck{\ill{}} non
\uncertain{unipotent} (i.e.\ formé d'élts \uncertain{semi-simples}, i.e.\
contenu dans un tore maximal $T$)

\[
  \{\Gamma\} \simeq \mathrm{or}(\Gamma_{0}) \simeq \text{gén}(Z)
  \simeq \text{gén}(Z) \simeq \mathrm{Inv}(T)
\]
\note{sous les termes, les numéros cerclés (1), (2), (3), (3'), (4) ; le mot
devant $\{\Gamma\}$ est illisible (\ill{}), et « $\mathrm{Inv}$ » est une
lecture douteuse. En marge droite, écrite dans le sens de la hauteur, la
suite de la chaîne : « $\simeq \mathrm{Fix}(u) \simeq$ \uncertain{$B_{S} -
(B_{S})^{Z}$} », avec les numéros (5) et (6)}

(3') Variété \uncertain{complexe} $V$ de rg 1 + (réseau $\Pi$) +
($Z \subset \Pi$) + torseur $X$ sous $V/\Pi$
\note{sous $Z \subset \Pi$, en petit : « facteur direct de rg 1 »}

\page{13}
\note{la face de l'enveloppe dont la page 12 est le dos : elle porte le
cachet de l'université de Nagoya daté du 20.12.77. Les calculs qui suivent
y sont écrits tête-bêche par rapport à l'adresse}

$\{\lambda, \lambda^{-1}\}$ \qquad $|\lambda| \neq 1$ \qquad donc
$\lambda \neq \lambda^{-1}$ i.e.\ $\lambda^{2} \neq 1$ i.e.\
$\lambda \neq \pm 1$

$\lambda + \lambda^{-1} = \alpha$ \quad tel que
$\alpha \notin [-2, +2]$ \qquad $\alpha^{2} - 4$

$\lambda\bar{\lambda} = 1$ \qquad $\bar{\lambda} = \lambda^{-1}$ \qquad
$\lambda + \lambda^{-1}$ réel

\section{Lettre à Lipman Bers (Les Aumettes, 15.4.1984)}
\note{lettre dactylographiée, en anglais, de six pages numérotées 1 à 6,
sans signature sur ce feuillet ; les corrections manuscrites sont intégrées
à leur place et signalées quand elles remplacent un mot tapé. Les lettres
grecques ($\nu$, $\lambda$, $\rho$), le $l$ des longueurs et plusieurs
indices sont portés à la main dans la frappe}

\page{15}
Les Aumettes 15.4.1984

Dear Lipman Bers,

Together with Yves Ladegaillerie (a former student of mine) we are running a
microseminar on the Teichmüller spaces and groups, my own motivations coming
mainly from algebraic geometry, and Ladegaillerie's from his interest in the
topology of surfaces. Lately we have met with a problem which I would like to
submit to you, as I understand you are the main expert on Thurston's
hyperbolic geometry approach to Teichmüller space. Before stating the specific
problem on hyperbolic “pants” (which things boil down to), let me tell you
what we are really after.

Assuming given a compact oriented surface with boundary $X_{0}$
\struck{\ill{}} as a reference-surface for constructing the Teichmüller-type
spaces, of genus $g$ and with “holes” (satisfying $2g-2+\nu > 0$), my primary
interest is in the more “algebraic” version of Teichmüller space,
corresponding to the question of classifying algebraic non singular curves
over $\underline{C}$, of genus $g$, with a system of \add{$\nu$} points (all
distinct) given on $X$, together with a “Teichmüller rigidification” of
$(X,S)$ namely a homotopy equivalence between $X_{0}$ and $X \setminus S$.
I'll denote this space, homeomorphic to $\underline{C}^{d}$ (where
$d = 3g-3+\nu$), by $\widetilde{M}_{g,\nu}$ (the tilda suggesting that it is
the universal covering of a finer object I am still more interested in,
namely the algebraic variety (or rather “multiplicity”, or “stack” in the
terminology of Mumford-Deligne) of moduli for algebraic curves of type
$(g,\nu)$). Thurston however considers a different modular space, where
algebraic curves with a given system of points are replaced by compact
conformal oriented surfaces \emph{with boundary}, giving rise to a modular
space $\widetilde{MB}_{g,\nu}$ (where the letter B recalls that we are
classifying structures with boundary) homeomorphic to
$\underline{C}^{d} \times (\underline{R}^{*+})^{\nu}$, where the extra factor
\struck{\ill{}} corresponds to the extra parameters introducing through the
existence of the boundary, namely the length's of the \struck{\ill{}}
components of the boundary with respect to the canonical hyperbolic structure
on the given surface. Our interest is in pinpointing the precise relationship
between the two modular spaces. The obvious idea here is to consider the case
of an algebraic curve with $\nu$ points given as a limit-case of a compact
conformal surface with boundary, when all the lengths $l_{i}$ of the
components of the boundary tend to zero. Therefore, it looks suitable to
consider both modular spaces above as imbedded in a larger third one, which
corresponds to the same modular problem as in Thurston's theory, except that
we allow the “boundary” to have some components reduced to just one point, in
the neighbourhood of which $X$ is just a conformal surface without boundary,
but with a given point
\note{la frappe laisse en blanc le nombre de points après « a system of » ;
« compact », devant « conformal surface with boundary », est ajouté à la main}

\page{16}
(viewed as a component of such a “generalised boundary”). We now should
\struck{should} get a modular space for “compact conformal oriented surfaces
with generalized boundary” (of type $g,\nu$ and rigidified via $X_{0}$), call
it $\widetilde{MB}_{g,\nu}$, homeomorphic to
$\underline{C}^{d} \times (\underline{R}^{+})^{\nu}$, where now the second
factor corresponds to the “parameters” $l_{i}$, which are allowed to take also
value 0 (which means that the corresponding component of the generalized
boundary is just one point). Thus $\widetilde{MB}_{g,\nu}$ appears as a
variety with boundary (in the topological sense — in the real analytic sense,
the “boundary” admits “corner-like” points obviously), and
$\widetilde{M}_{g,\nu}$ appears as a part of the boundary.
\note{le $+$ de $(\underline{R}^{+})^{\nu}$ est récrit à la main sur un
autre signe tapé}

My interest is in a better geometric understanding of the situation, which
should be “intrinsic” namely not depend on any particular choice of a
surgical decomposition of the reference surface $X_{0}$ into “pants”, used in
order to describe in a handy way \struck{the} standard “coordinate functions”
on the modular space $\widetilde{MB}_{g,\nu}$. There appears to be a
geometrically meaningful retraction of $\widetilde{MB}_{g,\nu}$ upon
$\widetilde{M}_{g,\nu}$ (commuting to the operations of the Teichmüller
modular group), the fibers being homeomorphic to
\uncertain{$(\underline{R}^{+})^{\nu}$} — more specifically, I expect the
semi-group $(\underline{R}^{+})^{I}$ (where $I$ is the set of indices for the
“holes” of $X_{0}$) to act on $MB$ in a natural way, with free action of the
subgroup $(R^{+*})^{\nu}$ upon $\widetilde{MB}^{o}$, in such a way that
$\widetilde{M}$ is just the quotient of $\widetilde{MB}$ by this action (or of
$\widetilde{MB}^{o}$ by the action of the \struck{\ill{}} corresponding
subgroup), and that \struck{each} fiber $F$ is isomorphic to
$(\underline{R}^{+})^{I}$ by the choice of any “origin” in
$F \cap \widetilde{MB}^{o}$.
\note{l'exposant de $(\underline{R}^{+})$ après « homeomorphic to » est
récrit à la main et se lit mal ; « each » tapé est corrigé à la main, et
« $F$ » ajouté}

Of course, “computainnally”, in terms of a decomposition of $X_{0}$ into
pants, the idea of such an operation is pretty obvious — namely letting the
components $\lambda_{i}$ of $\lambda \in (\underline{R}^{+})^{I}$ act as a
“multiplier” on the corresponding coordinate $l_{i}$. However, it is not
clear that this operation is intrinsic — and if it were intrinsic, an
intrinsic geometric description would still be \struck{\ill{}} desired.

Of course, in the description of the situation proposed above, the
retraction of $\widetilde{MB}$ upon $\widetilde{M}$ is obtained by
multiplying with the 0 multiplier (all $\lambda_{i}$ are 0). Now there is a
direct geometrical description of \struck{this} a retraction, by hyperbolic
surgery. Namely, for any compact conformal surface of type $g,\nu$ with
generalized boundary, let's “fill in” the holes which correspond to ordinary
components of the boundary, which are riemanian oriented circles, by “gluing
in” the cones on these circles (which are canonically endowed with a
conformal structure, using the riemanian structure on the given circles).
Thus we get a “functor” from compact conformal surfaces with generalized
boundary (of type $g,\nu$) to compact conformal surfaces \emph{without
boundary}, endowed with a system of $\nu$ points (making up a

\page{17}
“wholly degenerate” generalized boundary). When we throw in the
rigidifications and go over to isomorphism classes, this should give the
desired retraction. However, the geometric situation is a lot richer still,
as the compact surface without boundary obtained through surgery is endowed,
not only with a system of $\nu$ points, but moreover with a system of
mutually disjoint \emph{discs} around these points. The shape of these discs
is by no means arbitrary — we'll say that a system of discs around $\nu$
points on a compact conformal surface $X^{\wedge}$ without boundary is
“admissible”, if the situation can be obtained as above (up to isomorphism)
from surgery, starting with a compact conformal surface $X$ with boundary.
(NB Among the given “discs”, we should allow that some should be reduced to
their center — we'll call them “degenerate”.) The condition of admissibility
can be expressed intrinsically, by stating that \struck{the} for every
non-degenerate component $\Gamma_{i}$ of the system of boundaries of those
discs, the two operations we got of the standard circle group (of complex
numbers of module 1) upon $\Gamma_{i}$, by using the fact that it is (on the
one hand) the boundary of the disc $D_{i}$, and (on the other hand) that it
is a component of the boundary of the hyperbolic surface
$X^{\wedge} \setminus \bigl(\bigcup_{j} D_{j}^{o}\bigr)$, should be the same.
When $X^{\wedge}$ and the points $s_{i}$ on $X^{\wedge}$ are given, the
possible \struck{systems} admissible systems of discs around the points
$s_{i}$ depend on $\nu$ parameters — and the first idea which flips to mind
to give a more precise meaning to these “parameters”, is to view them as
being the “radii” of those discs. But then we'll have to define what we mean
by these !
\marginal{!}

The idea here is that, when we have a conformal disc $D$ and an interior
point $s$ of $D$, then $D$ may be viewed as canonically embedded in the
tangent space $T_{s}$ to $D$ at $s$, as the “unit disc” at $s$. Thus, in the
situation above of an \struck{\ill{}} admissible system of discs
$(D_{i})_{i\in I}$ around $(s_{i})_{i\in I}$, for every $s_{i}$ corresponding
to a non-degenerate $D_{i}$, we get a canonical disc
\[
  \Delta_{i} \subset T_{s_{i}}
\]
in the tangent space — and of course, for degenerate $D_{i}$, we'll take
$\Delta_{i}$ to be degenerate too. The discs we get in a given $T_{s_{i}}$
(for a fixed system $(s_{i})$, and a variable admissible system of discs
around these $s_{i}$) are all discs in the strict euclidean sense, given by
an unequality
\[
  |z| \leqslant r_{i} ,
\]
where $z \mapsto |z|$ denotes some hermitian metric on $T_{s_{i}}$ compatible
with the conformal structure — this metric being \struck{defined} unique up to
a scalar factor. the set $R_{i}$ of all those possible discs (the
non-degenerate ones, say) may be viewed in a natural way as a “torsor”
(= principal homogeneous space)
\note{« conformal » devant « surface $X^{\wedge}$ », « admissible » devant
« systems » (deux fois), « unique » et l'indice $R_{i}$ sont ajoutés à la
main}

\page{18}
under $\underline{R}^{+*}$, which plays here the role of the parameter space
of all possible (non degenerate) “radii” at $s_{i}$. If we admit also radius
zero, we accordingly get a parameter space $R_{i}^{\wedge}$, which may be
viewed as a torsor of sorts \struck{over} under $\underline{R}^{+}$. Thus the
set of radii for a given admissible system of discs $D_{i}$ around the points
$s_{i}$ may be viewed as a point of the product-space
\[
  r = (r_{i})_{i\in I} \in R^{\wedge} = \prod_{i\in I} R_{i}^{\wedge} .
\]
My expectation is that an admissible set of discs $(D_{i})$ is well
determined by the knowledge of the corresponding set of radii $r$, and
moreover that a given set of radii $r$ corresponds to an admissible system of
\struck{admissible} discs iff it satisfies a set of unequalities
\[
  r_{i} < \rho_{i} ,
\]
where
\[
  \rho = (\rho_{i})_{i\in I} \in R = \prod_{i\in I} R_{i}
\]
is \struck{a} some fixed system of radii, corresponding to a fixed system of
choices of hermitian metrics in the tangent spaces $T_{s_{i}}$.

I now see that this “expectation” does'nt quite match with the previous one,
about a “natural operation” of $(\underline{R}^{+})^{I}$ upon $\widetilde{MB}$,
having certain properties — it would match only if all $\rho_{i}$ where equal
to $+\infty$ (hence not in $R_{i}$ itself strictly speaking). I must confess I
did'nt look too thoroughly yet at the situation, and moreover I've been busy
with rather different kind of things for the last two or three months, and
lost a little contact \struck{contact}. What is clear however is that the
main key to an understanding of the general situation, is in an understanding
of the basic particular case of Thurston's pants. If we number $0, 1, \infty$
the three “holes” of such a pant, the surface $X^{\wedge}$ can be identified
canonically to the Riemann sphere, and the basic question then is to
understand how the pant is imbedded in this sphere $\Sigma$, as a complement
of \struck{(open)} the union of (open) discs around the points
$0, 1, \infty$, these discs forming an “admissible system”. So the main
question is about understanding the structure of all possible admissible
systems of three discs on $\Sigma$.

Puzzling a little about this problem, the following model came to my mind
(corresponding to “limiting radii” $\rho_{i}$ which are \emph{finite}, not
infinite). I view $\Sigma$ as endowed with it's usual euclidean metric, for
which the real projective line is a great circle, with $0, 1, \infty$ at
equal distance from each other on this equator. These points may be viewed
as the centers of three “orange \struck{quarters} slices”, making up a
cellular subdivision of $\Sigma$, where the common boundary of two among the
“\struck{quarters} slices” $Q_{i}$ ($i \in \{0, 1, \infty\}$) is a
half-great circle passing in between $s_{i}$ and $s_{j}$ at equal
\note{« all », « $r$ », « an admissible », « some », « with », « contact »
et les deux « slices » sont de sa main, en interligne ; dans
$(\underline{R}^{+})^{I}$, l'exposant $+$ est récrit sur un signe tapé}

\page{19}
distance from both, these three half-circles joining at the two poles $P^{+}$
and $P^{-}$. The “disc” $Q_{i}$ around $s_{i}$ has a conical structure around
$s_{i}$ (as has any conformal pointed disc), and we may take the concentric
\struck{discs} discs $\lambda_{i} Q_{i}$ with
\[
  0 < \lambda_{i} < 1 .
\]
The model I had in mind was that the (non degenerate) admissible systems of
discs around the points $s_{i}$ ($i \in \{0, 1, \infty\}$) \struck{\ill{}}
are exactly the systems of discs $\lambda_{i} Q_{i}$, with $\lambda_{i}$ as
above. (If we allow some discs to be degenerate, this means that instead of
the unequality above we merely demand $0 \leqslant \lambda_{i} < 1$, 0 not
excluded.)

This model, if correct, would give a rather precise description of the
inclusion relationships between pants, when these are considered as embedded
in the sphere. The intersection of all would be this system of three half
circles $C_{i}$, and the two poles $P^{+}, P^{-}$ would play a significant
role in the geometry of the pants, from this point of view. But it does'nt
seem that neither those half circles (which need not be geodesical I guess),
nor the two poles, have ever been described as intrinsically associated to a
pant. Of course, this model would give alternative “parameters” $\lambda_{i}$
for describing a pant, which are best suited for grasping the pants in terms
of spherical geometry. The next question would be an understanding of the
relationship between these parameters, and Thurston's $l_{i}$. Maybe it is
unreasonable to expect that for given index $i \in \{0, 1, \infty\}$, the
length $l_{i}$ depends only on $\lambda_{i}$ and not on the other parameters
$\lambda_{j}$ — and for this reason, \struck{\ill{}} the intuition at the
beginning of this letter, using Thurston's coordinate functions and notably
the $l_{i}$'s to get a fibration structure on $\widetilde{MB}$ over
$\widetilde{M}$, in terms of a given decomposition of $X_{0}$ into pants, is
probably not really relevant, namely it is non intrinsic. Assuming the model
I am suggesting is correct, the accurate description \struck{\ill{}} of
$\widetilde{MB}$ in terms of $\widetilde{M}$ would be
\[
  \widetilde{MB} \simeq \widetilde{M} \times [0, 1[^{I} ,
\]
where the second factor on the right hand sight refers to the system of
multipliers $\lambda_{i}$ ($i \in I$), tied to the $r_{i}$ above by
$r_{i} = \lambda_{i} \rho_{i}$.

My question of course is whether you have any information or idea to
propose, especially on the basic problem of relying pants to spherical
geometry, and more specifically, whether the model above is likely to hold,
or is definitely false. Also, one difficulty we found with hyperbolic
geometry of conformal surfaces, is that apart from existence and unicity of
the hyperbolic structure (compatible with the given conformal one and for
which the boundary is geodesic), there seems to be little

\page{20}
hold on more specific properties. As an example, starting with a compact
conformal surface with boundary $X$ (a pant, say), of hyperbolic type, and
removing an (open) “collar” around the boundary, we get another surface with
boundary $X'$ — what about the relation \struck{of} between the two
corresponding metrics ? Assuming the model for pants above is correct, it
would be nice to have an explicit expression of the metric of a pant in terms
of the parameters $\lambda_{i}$.

With my thanks for your attention, and for whatever comment you will care to
make, very sincerely ypur's
\note{« compact » et « between » sont ajoutés à la main ; la lettre s'arrête
sur la formule de politesse, sans signature sur ce feuillet}

\end{document}
